\documentclass[11pt]{article}
\usepackage[margin=2.4cm]{geometry}
\usepackage{fontspec}
\setmainfont{DejaVu Sans}
\usepackage{juliamono}
\usepackage{fvextra}
\usepackage{xcolor}
\usepackage{xurl}
\usepackage{hyperref}
\usepackage{amssymb}
\newcommand{\checkbox}{$\square$}
\newcommand{\answerchoices}{\checkbox\ YES \quad \checkbox\ NO \quad \checkbox\ CANNOT DETERMINE}
\newcommand{\decisionexplanation}{%
  \par\smallskip\textbf{If NO, give the reason for the mismatch (required). If
  CANNOT DETERMINE, state what context or evidence is missing:}\par
  \rule{\linewidth}{0.4pt}\par\vspace{7pt}\rule{\linewidth}{0.4pt}\par}
\setlength{\parindent}{0pt}\setlength{\parskip}{5pt}
\setlength{\emergencystretch}{2em}
\DefineVerbatimEnvironment{Vsrc}{Verbatim}{fontsize=\small,frame=single,
  breaklines=true,breakanywhere=true}
\DefineVerbatimEnvironment{Vlean}{Verbatim}{fontsize=\scriptsize,frame=single,
  numbers=left,numbersep=5pt,breaklines=true,breakanywhere=true}
\hypersetup{colorlinks=true,urlcolor=blue}

\begin{document}
\begin{center}{\LARGE Proof Verification Document}\\[2pt]{\large mse\_811862\_\_ans811879}\\[2pt]{\small math.StackExchange curated corpus (confirmed)}\end{center}
\begin{center}\fcolorbox{blue!60!black}{blue!4}{\parbox{0.92\linewidth}{\textbf{Make three independent decisions.} Follow the evidence chain: original source $\longrightarrow$ challenged claim $\longrightarrow$ proposed counterexample $\longrightarrow$ Lean theorem. A record is valid only when all three links are correct. Complete the human sign-off before reading the automated review appendix; those earlier results are background, not an answer key.}}\end{center}
\textbf{Use only this document, except for the public source link.} Everything needed for Questions 2 and 3 is included here. For Question 1, open the public source link when one is shown. If there is no public link, the complete stored problem and solution are included here. If the information is still not enough, choose CANNOT DETERMINE.
For each question, choose \textbf{YES}, \textbf{NO}, or \textbf{CANNOT DETERMINE}. Choose CANNOT DETERMINE if the document does not provide enough information for you to decide; do not mark YES if you are uncertain. If you find a specific mismatch, choose NO and explain it in the space below that question. Do not rewrite the claim, add assumptions, or substitute a different counterexample.
\textbf{How copied text is shown.} The problem, quoted passage, and proposed counterexample come from the stored record. The Lean theorem statement and Lean files are copied exactly from the certificate. Some long lines continue on the next printed line; this is only page formatting and does not change the copied text. Line numbers are added only to help you find passages. If the stored record contains a character that cannot be printed, the document shows \texttt{[U+XXXX CONTROL]} in its place.
\begin{description}
\item[Challenged claim.] The exact statement from the source that this record says is false.
\item[Proposed counterexample.] A concrete case that should satisfy the claim's assumptions but make its conclusion false.
\item[Lean certificate.] A Lean file proving a formal theorem. Successful compilation proves the theorem as written, not that it is the right translation of the source.
\end{description}
\section*{Question 1 --- Does this proof verification document accurately represent the source?}
\textbf{What to decide.} Compare the challenged passage below with the original source---either the public link or the complete source text in this document. Answer YES only if they make the same claim and include the same assumptions or conditions. If this document changes the meaning, leaves out a condition, or points to a different passage, answer NO.
\subsection*{Original source}
\textbf{Source record:}
\begin{Vsrc}
qedbench_certificate_availability_triage/data/records_mse.jsonl
\end{Vsrc}
\textbf{Public source:}\par
\url{https://math.stackexchange.com/a/811879}\par
\subsection*{Problem or question from the source (verbatim)}
\begin{Vsrc}
Title: Gaussian Matrix Integral

I need your help to solve this exercise :

Let $S$ be a symmetric Hermitian matrix $N\times N$ : $S=(s_{ij})$ with $s_{ij}=s_{ji}$.

When $\langle s_{ij}s_{kl}\rangle\neq 0$
What $$\int Tr(S^{2n})\;d\mu(s)$$ counts?
Calculate the moments of the matrix trace given by $$
\int Tr(S^2) dS,\quad \int Tr(S^4) dS
$$
where S is a real symmetric matrix satisfying $S_{ij}=S_{ji}$.  

We are supposed to use The Wick Theorem to solve it. 

any idea on how to calculate this, thank you.
\end{Vsrc}
\subsection*{Challenged passage recorded from the source (stored text)}
\begin{Vsrc}
\langle Tr(H^{2k})\rangle= \sum_{i_1,i_2,...i_{2k}}H_{i_1 i_2} H_{i_2 i_3}...H_{i_{2k}i_1}=\frac{1}{k+1}\binom {2k} {k}=C_k
\end{Vsrc}
\subsection*{Surrounding source lines (stored text with added line labels)}
The labels refer to lines in the solution text stored with the source record; the source-line content is not transliterated.
\begin{Vsrc}
0059 | \langle Tr(H^{4})\rangle=\langle H_{ij}H_{jk}\rangle \langle H_{kl}H_{li}\rangle + \langle H_{ij} H_{kl}\rangle \langle H_{jk}H_{li}\rangle + \langle H_{ij}H_{li}\rangle \langle H_{jk} H_{kl} \rangle 
0060 | \end{equation}
0061 | where a contraction is defined by 
0062 | \begin{equation}
0063 | \langle H_{ij}H_{kl}\rangle=\delta_{il} \delta_{jk}.
0064 | \end{equation}
0065 | Using this result and contracting indices we can write 
0066 | \begin{equation}
0067 | \langle Tr(H^{4})\rangle=\sum_{i,j,k,l} (\delta_{ik}\delta_{jj}\delta_{ki}\delta_{ll}+\delta_{il}\delta_{kj}\delta_{ji}\delta_{lk}+\delta_{ii}\delta_{jl}\delta_{lj}\delta_{kk})=2N^3+N,\qquad (\text{answer  to  c})
0068 | \end{equation}
0069 | since one of the contractions only has one free index it contributes N.  The other two contractions both contribute N$^3$.  The 6th moment of the distribution has 15 possibilities to contract so I will omit this derivation.  In general, the nth moment is written by expanding the product as
0070 | \begin{equation}
0071 | \langle Tr(H^{2k})\rangle= \sum_{i_1,i_2,...i_{2k}}H_{i_1 i_2} H_{i_2 i_3}...H_{i_{2k}i_1}=\frac{1}{k+1}\binom {2k} {k}=C_k \qquad (\text{general answer to A,B,C})
0072 | \end{equation}
0073 | where C$_{k}$ is the Catalan numbers which represents the k different ways to match the entries.  This result can be plugged into the average resolvent expression and summed to give
0074 | \begin{equation}
0075 | \langle G(z)\rangle = \frac{1}{N} \sum_{k=0}^{\infty} z^{-(2k+1)}\langle 
0076 | Tr(H^{2k})\rangle=\frac{1}{2}\big(z-\sqrt{z^2+4}\big).
0077 | \end{equation}
0078 | This equality is obtained by using contractions and expanding the average resolvent for large z as done above \begin{equation}
0079 | \langle G(z) \rangle=\bigg\langle \frac{1}{N} \big(z^{-1}+z^{-2}TrH+z^{-3}Tr H^2 +z^{-4} Tr H^3 + O(N)\big )\bigg \rangle.
0080 | \end{equation}
0081 | Since the odd moments vanish due to symmetry, the even moments are calculated as shown above.  We obtain
0082 | \begin{equation}
0083 | \langle G(z) \rangle= z^{-1}+z^{-1}\langle G(z) \rangle ^2
\end{Vsrc}
\textbf{A YES means:} the proof verification document faithfully represents both the source context and the challenged claim, without dropping a condition or changing its meaning.
\section*{Question 2 --- Does the proposed counterexample disprove that claim?}
\textbf{Claim to test (repeated for convenience).}
\begin{Vsrc}
\langle Tr(H^{2k})\rangle= \sum_{i_1,i_2,...i_{2k}}H_{i_1 i_2} H_{i_2 i_3}...H_{i_{2k}i_1}=\frac{1}{k+1}\binom {2k} {k}=C_k
\end{Vsrc}
\subsection*{How the record says the argument uses this claim}
This is the record's description, not an established fact; check it against the source.
\begin{Vsrc}
The proof's stated general formula for the trace moments, $\langle Tr(H^{2k})\rangle = \frac{1}{k+1}\binom{2k}{k} = C_k$ (labelled "general answer to A,B,C"), asserts that the 2k-th trace moment equals the Catalan number $C_k$ with no dependence on the matrix size $N$. This contradicts the proof's own explicitly computed values $\langle Tr(H^2)\rangle = N^2$ and $\langle Tr(H^4)\rangle = 2N^3 + N$, and contradicts the proof's own Wick contraction rule $\langle H_{ij}H_{kl}\rangle=\delta_{il}\delta_{jk}$. (The correct statement is that $C_k$ is the leading large-$N$ coefficient, i.e. $\langle Tr(H^{2k})\rangle \sim C_k N^{k+1}$, not the moment itself.)
\end{Vsrc}
\subsection*{Proposed counterexample (exact values or construction)}
\begin{Vsrc}
k = 1, N = 2 (smallest witness; 29 of the 30 exhaustively tested (k,N) pairs falsify the claim). Using the proof's own evaluation rule <H_ij H_kl> = delta_il delta_jk together with Gaussian/Wick pairing, the exactly computed integer is <Tr(H^2)> = sum_{i,j} <H_ij H_ji> = sum_{i,j} delta_ii delta_jj = N^2 = 4, whereas the asserted value is C_1 = binom(2,1)/2 = 1. 4 != 1. Further exact witnesses from the same exhaustive table: (k=2, N=2): <Tr H^4> = 18 vs C_2 = 2; (k=3, N=2): <Tr H^6> = 120 vs C_3 = 5; (k=4, N=3): <Tr H^8> = 5355 vs C_4 = 14; (k=5, N=6): <Tr H^10> = 2521260 vs C_5 = 42. The only pair where the asserted equality holds is the degenerate (k=1, N=1).
\end{Vsrc}
\subsection*{Why the record says the counterexample works}
Check this explanation independently; it is evidence to assess, not an answer key.
\begin{Vsrc}
The sentence asserts a literal equality between two completely determined integers for each (k, N): the left side is the Gaussian expectation of Tr(H^{2k}), which the proof itself prescribes how to evaluate (centred Gaussian entries plus the propagator <H_ij H_kl> = delta_il delta_jk, hence Wick/Isserlis: sum over all (2k-1)!! perfect matchings of the 2k factors of the product of propagators); the right side is the Catalan number C_k = binom(2k,k)/(k+1). Nothing asymptotic or infinitary is stated — no limit, no ~, no O(-) — so a single (k, N) pair with unequal integers refutes it, and the claim is finitely refutable.

The script computes the left side exactly, in Python arbitrary-precision integers only (no float appears in any decision path), by two independent methods that are cross-checked for agreement wherever affordable:
  METHOD A: for each perfect matching of the 2k factors, the deltas identify summation index variables i_1..i_{2k}; union-find gives the number c of connected components, and the number of admissible tuples in {1..N}^{2k} is exactly N^c. Contribution = N**c (exact integer power); sum over all matchings.
  METHOD B: literal enumeration of all N^{2k} index tuples against all matchings, counting exact delta satisfaction. Run for every (k,N) pair whose cost N^{2k} * (2k-1)!! <= 3e6 (this covers all of k<=3, plus k=4 with N<=3 and k=5 with N=2). Method B agreed with Method A on every pair where it ran; a mismatch would have aborted the script.

Validation that the implementation encodes the proof's own model: the exact values produced reproduce, for every N in 1..6, precisely the two moments the proof itself computes explicitly — <Tr H^2> = N^2 and <Tr H^4> = 2N^3 + N (e.g. N=3 gives 9 and 57; N=6 gives 36 and 438). So the refutation is internal to the proof: the same rule that yields the proof's own N^2 and 2N^3+N cannot also yield the N-independent C_k, and the general formula is inconsistent with the proof's own worked cases. The correct statement is the asymptotic one, <Tr(H^{2k})> = C_k N^{k+1} + (lower order), which the printed column C_k*N^{k+1} illustrates: e.g. k=5, N=6 gives exact 2521260 against leading term 1959552 — close in relative order but never equal, and both wildly different from C_5 = 42.

Runtime 0.85 s. The claim fails on 29 of the 30 exhaustively scanned pairs; the sole exception is the degenerate N=1, k=1.
\end{Vsrc}
\subsection*{How to decide}
\begin{enumerate}
\item List every assumption that applies to the challenged claim.
\item Check that the proposed example satisfies every one of those assumptions.
\item Recompute the claimed outcome and the actual outcome using exact reasoning.
\item Confirm that the conflict is with the challenged claim itself, not an unrelated sentence.
\end{enumerate}
\textbf{A YES means:} the example is fully specified, satisfies all relevant assumptions, and makes the challenged claim false. If a necessary value, condition, or calculation is missing, choose CANNOT DETERMINE.
\section*{Question 3 --- Does the final Lean theorem statement verify this same counterexample?}
\textbf{What to decide.} A Lean file can compile while proving the wrong theorem. Compare these three links independently:
\begin{enumerate}
\item The Lean definitions represent the same mathematical objects, assumptions, and operations as the challenged claim.
\item The Lean construction uses the same concrete values or example shown in Question 2.
\item The final theorem statement asserts that the challenged claim fails on that example, not merely a weaker or nearby statement.
\end{enumerate}
\subsection*{Selected final Lean theorem statement (verbatim UTF-8 excerpt)}
This is an exact contiguous slice of the certificate file, from the theorem keyword through the original \texttt{:=} token where its proof begins. It is shown first so the formal target is not buried in the complete file.
\begin{Vlean}
theorem claim_fails_1_2 : claimHolds 1 2 = false :=
\end{Vlean}
\subsection*{Project-local imported Lean files (complete, verbatim UTF-8)}
This section follows the certificate's import commands recursively and prints every repository-owned Lean source file on which the final theorem statement depends. Each file is copied in full, including its definitions and controls. External Mathlib, Std, and Lean Init source is not duplicated.
No project-local imported Lean files were found. External Mathlib, Std, and Lean Init source is not duplicated in this proof verification document.
\subsection*{Complete counterexample Lean file (verbatim UTF-8)}
The complete Lean file is included below so you can check every definition and supporting result used by the final theorem statement. Use it together with the project-local imported Lean files above.
\textbf{Lean certificate file:}
\begin{Vsrc}
partB_feasibility/mse_certificates/mse_811862__ans811879.lean
\end{Vsrc}
\begin{Vlean}
/-!
# mse_811862__ans811879 — self-contained refutation certificate

* Record id : `mse_811862__ans811879`
* Question  : https://math.stackexchange.com/questions/811862  (inferred from the record id)
* Answer    : https://math.stackexchange.com/a/811879            (inferred from the record id)
* Route     : SELF_CONTAINED_LEAN (no `import`, no Mathlib; kernel `decide` only)
* Text licence: CC BY-SA (Stack Exchange).

## The claim being refuted

The refuted sentence asserts a LITERAL EQUALITY, for each `k ≥ 1` and each matrix size
`N ≥ 1`, between the Gaussian (GUE-style) expectation of the trace of the `2k`-th power
of an `N × N` random matrix `H` and the `k`-th Catalan number:

>   ⟨Tr(H^{2k})⟩  =  C_k  =  binom(2k, k) / (k+1).

Nothing asymptotic or infinitary is stated in that sentence — no limit `N → ∞`, no `~`,
no normalisation by a power of `N`, no `O(·)`.  Both sides are, for each fixed `(k, N)`,
completely determined integers:

* the RIGHT side is the Catalan number `C_k`, an `N`-independent integer;
* the LEFT side is fixed by the proof's OWN evaluation rule.  The proof declares the
  entries centred Gaussian with propagator

      ⟨H_ij H_kl⟩ = δ_il δ_jk,

  and therefore evaluates any even moment by Wick/Isserlis: sum over all `(2k-1)!!`
  perfect matchings of the `2k` factors of the product of the propagators.

So a single pair `(k, N)` with unequal integers refutes the sentence, and the claim is
finitely refutable.  (The CORRECT statement is the asymptotic one,
`⟨Tr(H^{2k})⟩ = C_k · N^{k+1} + lower order`; the asserted `N`-free equality is not it.)

## Witness

    k = 1, N = 2 :   ⟨Tr(H^2)⟩ = 4   vs   C_1 = binom(2,1)/2 = 1.        4 ≠ 1.

This is the SMALLEST witness.  Expanding it by hand from the proof's own rule:

    ⟨Tr(H^2)⟩ = Σ_{i,j} ⟨H_ij H_ji⟩ = Σ_{i,j} δ_ii δ_jj = N^2 = 4.

Further exact witnesses, all recomputed by the kernel below:

    (k=2, N=2) :  ⟨Tr H^4⟩  = 18       vs  C_2 = 2
    (k=3, N=2) :  ⟨Tr H^6⟩  = 120      vs  C_3 = 5
    (k=4, N=3) :  ⟨Tr H^8⟩  = 5355     vs  C_4 = 14
    (k=5, N=6) :  ⟨Tr H^10⟩ = 2521260  vs  C_5 = 42      (leading term C_5·N^6 = 1959552)

Of the 30 exhaustively scanned pairs (`1 ≤ k ≤ 5`, small `N`), 29 falsify the claim.  The
SOLE pair on which the asserted equality holds is the degenerate `(k, N) = (1, 1)`, where
`⟨Tr H^2⟩ = 1^2 = 1 = C_1` by coincidence.  That exception is not swept under the rug: it
is asserted positively in `scan_k3_N4` and in `nc2_sole_agreement_is_1_1` below, so the
refutation predicate is demonstrably not a predicate that fires on everything.

## How the left side is computed here (exactly, in `Nat`, two independent ways)

Write the trace as `Tr(H^{2k}) = Σ_{i_0 … i_{2k-1}} H_{i_0 i_1} H_{i_1 i_2} ⋯ H_{i_{2k-1} i_0}`,
so factor `t` (for `t = 0, …, 2k-1`) is `H_{a_t b_t}` with `a_t = t`, `b_t = t+1 mod 2k`.
Under the proof's propagator, pairing factor `s` with factor `t` contributes
`δ_{a_s b_t} δ_{b_s a_t}`, i.e. it FORCES the two index identifications

    i_{a_s} = i_{b_t}      and      i_{b_s} = i_{a_t}
    (that is:  s ~ t+1 (mod 2k)   and   s+1 ~ t (mod 2k) ).

**METHOD A** (`wickTrace`, the primary definition).  For each perfect matching of the
`2k` factors, impose those `2·k` identifications on the `2k` index VARIABLES by
union–find, count the number `c` of resulting equivalence classes, and add `N^c` — the
exact number of index tuples in a set of size `N` satisfying every delta.  Sum over all
`(2k-1)!!` matchings.

**METHOD B** (`wickTraceBrute`, an independent recomputation).  Literally enumerate all
`N^{2k}` index tuples and, for each matching, count by hand how many satisfy every delta.
Method B is exponentially more expensive, so it is run only on the affordable pairs; the
file PROVES the two methods agree on each of those (`nc7_methods_agree`).  A slip in the
union–find/component count could therefore not by itself produce the refutation.

All arithmetic is `Nat` (arbitrary precision, exact).  No floats, no rationals, no
division appears anywhere in the computation of the left side.  The only division in the
file is the exact `binom(2k,k) / (k+1)` inside `catalan`, and that value is independently
cross-checked against `binom(2k,k) - binom(2k,k+1)` (`nc4_catalan_two_ways`).

## Why this is INTERNAL to the proof, not an outside standard imposed on it

The very same evaluator, run at `k = 1` and `k = 2`, reproduces the two moments the proof
itself computes explicitly and by hand:

    ⟨Tr H^2⟩ = N^2        and       ⟨Tr H^4⟩ = 2N^3 + N

for every `N = 1, …, 6` (e.g. `N = 3` gives 9 and 57; `N = 6` gives 36 and 438).  This is
proved as `nc1_reproduces_proofs_own_moments`.  Hence the refutation is not a dispute
about which random-matrix model is meant: the rule that yields the proof's own `N^2` and
`2N^3 + N` cannot simultaneously yield the `N`-independent `C_k`, and the proof's general
formula is inconsistent with the proof's own worked cases.

## Modelling judgments (the kernel checks the arithmetic, not the faithfulness of these)

1. **Propagator.**  `⟨H_ij H_kl⟩ = δ_il δ_jk` is taken verbatim from the proof (it is the
   GUE/Wigner normalisation with unit variance and no `1/N`).  Any overall variance
   rescaling `σ²` multiplies every `2k`-moment by the SAME constant `σ^{2k}`, so it cannot
   turn an `N`-dependent left side into the `N`-free `C_k`; the refutation is robust to
   that choice.  What it is NOT robust to is silently inserting the `1/N^{k+1}` (or
   `1/N` per entry) normalisation that the correct asymptotic statement carries — but the
   refuted sentence contains no such normalisation, which is precisely the error.
2. **Centred + Gaussian ⇒ Wick.**  Odd moments vanish and even moments are the sum over
   perfect matchings of products of pairwise covariances (Isserlis' theorem).  This is the
   proof's own stated evaluation method; the file implements it, it does not assume it.
3. **Index alphabet.**  Indices range over a set of size `N`; the file uses `{0,…,N-1}`.
   Only equalities between indices ever matter, so the labelling is immaterial.
4. **Union–find is by global relabelling** (`unite` rewrites every occurrence of the
   larger label), so after a single pass over the constraint list the labels are exactly
   the equivalence classes of the transitive closure.  `nc5_components_have_teeth` pins
   the component counter down on hand-checked cases (including a matching that collapses
   all four variables into ONE class), so it is not returning a constant.
5. **Matching enumeration is complete and duplicate-free** by the standard "pair the first
   element with each other element, recurse on the rest" scheme, run with explicit fuel
   (structural recursion, kernel-reducible).  Its cardinality is verified to be the double
   factorial `(2k-1)!! = 1, 3, 15, 105, 945` (`nc6_matching_counts`); an under- or
   over-counting enumeration would fail that test.
6. **Fail-closed fuel.**  `matchingsAux` with fuel `0` returns the EMPTY list of matchings,
   not `[[]]`; an under-fuelled call therefore yields the value `0`, which agrees with no
   Catalan number `C_k ≥ 1` for the wrong reason — but `nc6_matching_counts` and
   `nc7_methods_agree` both fail loudly in that case, so exhausted fuel cannot silently
   manufacture a refutation.  (`matchingsOf n` is called with fuel `n+1`, ample.)
7. **The claim's own conclusion is expressible and can be TRUE in this model** — it holds
   at `(k,N) = (1,1)` (`nc2_sole_agreement_is_1_1`).  `claimHolds` is not false by
   construction.
8. Everything is discharged by kernel `decide`.  No `native_decide`, no `axiom`, no
   `sorry`, no well-founded recursion, no `partial`.
-/

set_option maxRecDepth 100000
set_option maxHeartbeats 4000000

namespace MSE_811862

/-! ## §0. Tiny list utilities (self-contained; no core `List` API relied upon) -/

def lmap {α β : Type} (f : α → β) : List α → List β
  | [] => []
  | a :: r => f a :: lmap f r

def lapp {α : Type} : List α → List α → List α
  | [], m => m
  | a :: r, m => a :: lapp r m

def joinL {α : Type} : List (List α) → List α
  | [] => []
  | l :: ls => lapp l (joinL ls)

def llen {α : Type} : List α → Nat
  | [] => 0
  | _ :: r => 1 + llen r

def allB {α : Type} (p : α → Bool) : List α → Bool
  | [] => true
  | a :: r => p a && allB p r

def sumN : List Nat → Nat
  | [] => 0
  | a :: r => a + sumN r

/-- `[0, 1, …, n-1]`. -/
def upto : Nat → List Nat
  | 0 => []
  | k + 1 => lapp (upto k) [k]

/-- `i`-th entry of a `Nat` list; `0` out of range.  All indices used below come from
`upto m` applied to lists of length exactly `m`, so the default is never reached. -/
def getD : List Nat → Nat → Nat
  | [], _ => 0
  | a :: _, 0 => a
  | _ :: r, k + 1 => getD r k

def memN (x : Nat) : List Nat → Bool
  | [] => false
  | a :: r => Nat.beq a x || memN x r

/-- Delete the first occurrence of `x`. -/
def removeFirst (x : Nat) : List Nat → List Nat
  | [] => []
  | a :: r => match Nat.beq a x with
    | true => r
    | false => a :: removeFirst x r

/-- Exact `b^e`, structurally recursive (kernel-reducible; `Nat.mul` is exact). -/
def npow (b : Nat) : Nat → Nat
  | 0 => 1
  | e + 1 => b * npow b e

/-- Number of distinct entries of a `Nat` list. -/
def distinctCount : List Nat → Nat
  | [] => 0
  | a :: r => (match memN a r with | true => 0 | false => 1) + distinctCount r

/-! ## §1. Perfect matchings of the `2k` matrix factors

Wick/Isserlis pairs up the `2k` factors of `Tr(H^{2k})` in every possible way.
`matchingsAux` is the textbook enumeration: pair the first element with each of the
others, recurse on the remainder.  It is complete and duplicate-free, and the fuel makes
it structurally recursive.  Fuel `0` returns NO matchings (fail-closed). -/

def matchingsAux : Nat → List Nat → List (List (Nat × Nat))
  | 0, _ => []
  | _ + 1, [] => [[]]
  | f + 1, x :: rest =>
      joinL (lmap (fun y =>
          lmap (fun m => (x, y) :: m) (matchingsAux f (removeFirst y rest)))
        rest)

/-- All perfect matchings of `{0, …, n-1}` (empty list when `n` is odd). -/
def matchingsOf (n : Nat) : List (List (Nat × Nat)) := matchingsAux (n + 1) (upto n)

/-! ## §2. The delta constraints imposed by one pairing

Factor `t` of `Tr(H^{2k})` is `H_{a_t b_t}` with `a_t = t`, `b_t = t+1 mod 2k`.
The proof's propagator `⟨H_ij H_kl⟩ = δ_il δ_jk` turns the pairing `{s, t}` into the two
index identifications `i_{a_s} = i_{b_t}` and `i_{b_s} = i_{a_t}`. -/

/-- Cyclic successor modulo `m` (`m = 2k`). -/
def nextIdx (m i : Nat) : Nat := (i + 1) % m

/-- The delta constraints (as pairs of index-variable names) of a whole matching. -/
def consOf (m : Nat) (mt : List (Nat × Nat)) : List (Nat × Nat) :=
  joinL (lmap (fun (p : Nat × Nat) => [(p.1, nextIdx m p.2), (nextIdx m p.1, p.2)]) mt)

/-! ## §3. METHOD A — union–find on the index variables, contribution `N^c`

`unite` merges by REWRITING every occurrence of the larger label with the smaller, so a
single pass over the constraint list leaves the labels equal to the equivalence classes
of the transitive closure of the constraints.  The number of index tuples over an
alphabet of size `N` satisfying all deltas is then exactly `N^{#classes}`. -/

def relabel (old new : Nat) (lbl : List Nat) : List Nat :=
  lmap (fun c => match Nat.beq c old with | true => new | false => c) lbl

def unite (lbl : List Nat) (x y : Nat) : List Nat :=
  let a := getD lbl x
  let b := getD lbl y
  match Nat.ble a b with
  | true => relabel b a lbl
  | false => relabel a b lbl

def uniteAll : List Nat → List (Nat × Nat) → List Nat
  | lbl, [] => lbl
  | lbl, p :: r => uniteAll (unite lbl p.1 p.2) r

/-- Number of independent index variables surviving one pairing. -/
def compOfMatching (m : Nat) (mt : List (Nat × Nat)) : Nat :=
  distinctCount (uniteAll (upto m) (consOf m mt))

/-- **`⟨Tr(H^{2k})⟩` for an `N × N` matrix, computed exactly (METHOD A).** -/
def wickTrace (k N : Nat) : Nat :=
  sumN (lmap (fun mt => npow N (compOfMatching (2 * k) mt)) (matchingsOf (2 * k)))

/-! ## §4. METHOD B — literal enumeration of all `N^{2k}` index tuples

Independent recomputation of the same quantity with no union–find and no exponentiation:
for every matching, walk the whole index space and count the tuples that satisfy every
delta.  Used only as a cross-check on affordable `(k, N)`. -/

def tuplesOf (N : Nat) : Nat → List (List Nat)
  | 0 => [[]]
  | len + 1 => joinL (lmap (fun t => lmap (fun v => v :: t) (upto N)) (tuplesOf N len))

def satisfies (v : List Nat) (cs : List (Nat × Nat)) : Bool :=
  allB (fun (p : Nat × Nat) => Nat.beq (getD v p.1) (getD v p.2)) cs

/-- **`⟨Tr(H^{2k})⟩` again (METHOD B): brute-force over the index space.** -/
def wickTraceBrute (k N : Nat) : Nat :=
  sumN (lmap (fun mt =>
      sumN (lmap (fun v => match satisfies v (consOf (2 * k) mt) with
              | true => 1
              | false => 0)
        (tuplesOf N (2 * k))))
    (matchingsOf (2 * k)))

def bruteAgrees (k N : Nat) : Bool := Nat.beq (wickTrace k N) (wickTraceBrute k N)

/-! ## §5. The right-hand side: Catalan numbers -/

/-- Pascal's rule; exact, no division. -/
def binom : Nat → Nat → Nat
  | _, 0 => 1
  | 0, _ + 1 => 0
  | n + 1, r + 1 => binom n r + binom n (r + 1)

/-- `C_k = binom(2k, k) / (k+1)` — the claim's own right-hand side, verbatim.
The division is exact on these arguments (checked against `catalanAlt`). -/
def catalan (k : Nat) : Nat := binom (2 * k) k / (k + 1)

/-- Independent formula `C_k = binom(2k, k) - binom(2k, k+1)`, used only to certify that
the division in `catalan` is the exact one intended. -/
def catalanAlt (k : Nat) : Nat := binom (2 * k) k - binom (2 * k) (k + 1)

/-! ## §6. The claim, and the refutation predicate -/

/-- The refuted sentence, instantiated at `(k, N)`: `⟨Tr(H^{2k})⟩ = C_k`. -/
def claimHolds (k N : Nat) : Bool := Nat.beq (wickTrace k N) (catalan k)

/-- **Refutation predicate.**  True exactly when `(k, N)` is an admissible instance of the
claim (`k ≥ 1`, `N ≥ 1`) and the asserted equality FAILS there. -/
def refutesAt (k N : Nat) : Bool :=
  Nat.ble 1 k && Nat.ble 1 N && !(claimHolds k N)

/-! ## §7. THE WITNESS — `k = 1`, `N = 2` -/

/-- The left side is exactly 4: `⟨Tr(H^2)⟩ = Σ_{i,j} δ_ii δ_jj = N² = 4`. -/
theorem lhs_1_2 : Nat.beq (wickTrace 1 2) 4 = true := by decide

/-- The right side is exactly 1: `C_1 = binom(2,1)/2 = 1`. -/
theorem rhs_1 : Nat.beq (catalan 1) 1 = true := by decide

/-- The asserted equality fails at the witness. -/
theorem claim_fails_1_2 : claimHolds 1 2 = false := by decide

/-- **THE REFUTATION.** -/
theorem refuted_1_2 : refutesAt 1 2 = true := by decide

/-- The same refutation with both integers spelled out in one statement. -/
theorem refuted_1_2_explicit :
    (Nat.beq (wickTrace 1 2) 4 && Nat.beq (catalan 1) 1 && refutesAt 1 2) = true := by
  decide

/-! ## §8. Further exact witnesses from the same exhaustive scan -/

/-- `(k=2, N=2)`: `⟨Tr H^4⟩ = 18` against `C_2 = 2`. -/
theorem refuted_2_2 :
    (Nat.beq (wickTrace 2 2) 18 && Nat.beq (catalan 2) 2 && refutesAt 2 2) = true := by
  decide

/-- `(k=3, N=2)`: `⟨Tr H^6⟩ = 120` against `C_3 = 5`. -/
theorem refuted_3_2 :
    (Nat.beq (wickTrace 3 2) 120 && Nat.beq (catalan 3) 5 && refutesAt 3 2) = true := by
  decide

/-- `(k=4, N=3)`: `⟨Tr H^8⟩ = 5355` against `C_4 = 14`. -/
theorem refuted_4_3 :
    (Nat.beq (wickTrace 4 3) 5355 && Nat.beq (catalan 4) 14 && refutesAt 4 3) = true := by
  decide

/-- `(k=5, N=6)`: `⟨Tr H^10⟩ = 2521260` against `C_5 = 42`.  For scale, the leading term
of the CORRECT asymptotic statement is `C_5 · N^6 = 42 · 46656 = 1959552` — the same order
of magnitude as the exact value, and nothing like `42`.  (Heaviest kernel computation in
the file: 945 pairings on 10 factors.) -/
theorem refuted_5_6 :
    (Nat.beq (wickTrace 5 6) 2521260
      && Nat.beq (catalan 5) 42
      && Nat.beq (42 * npow 6 6) 1959552
      && refutesAt 5 6) = true := by
  decide

/-! ## §9. Exhaustive scan

Over `1 ≤ k ≤ 3` and `1 ≤ N ≤ 4`: every pair refutes the claim EXCEPT the degenerate
`(1,1)`, where the claim is asserted to HOLD.  A single statement therefore certifies both
halves of the record's finding. -/

def scanRow (k : Nat) (Ns : List Nat) : Bool :=
  allB (fun N =>
      match Nat.beq k 1 && Nat.beq N 1 with
      | true => claimHolds k N          -- the sole coincidence, asserted positively
      | false => refutesAt k N)
    Ns

theorem scan_k3_N4 : allB (fun k => scanRow k [1, 2, 3, 4]) [1, 2, 3] = true := by decide

/-! ## §10. NEGATIVE CONTROLS

If the modelling were vacuous — a Wick evaluator that returned a constant, a matching
enumerator that produced nothing (or everything), a component counter stuck at `2k`, a
`claimHolds` false by construction, or a Catalan function that was not Catalan — the
theorems below would fail and the file would not compile. -/

/-- **NC1 — THE DECISIVE CONTROL: the evaluator reproduces the PROOF'S OWN worked cases.**
The proof computes, by hand and explicitly, `⟨Tr H^2⟩ = N^2` and `⟨Tr H^4⟩ = 2N^3 + N`.
The very same code that produces the refutation reproduces both, for every `N = 1,…,6`
(e.g. `N=3`: 9 and 57; `N=6`: 36 and 438).  So the refutation is internal to the proof:
the rule that gives the proof its own `N`-dependent answers cannot also give the
`N`-independent `C_k`. -/
def reproducesProofOwnMoments (N : Nat) : Bool :=
  Nat.beq (wickTrace 1 N) (N * N) && Nat.beq (wickTrace 2 N) (2 * N * N * N + N)

theorem nc1_reproduces_proofs_own_moments :
    allB reproducesProofOwnMoments [1, 2, 3, 4, 5, 6] = true := by decide

/-- **NC2 — the claim CAN be true in this model.**  At the degenerate `(k,N) = (1,1)` the
asserted equality genuinely holds (`⟨Tr H^2⟩ = 1 = C_1`), so `claimHolds` is not a
predicate that is false by construction and `refutesAt` does not fire on everything. -/
theorem nc2_sole_agreement_is_1_1 :
    (claimHolds 1 1 && Nat.beq (wickTrace 1 1) 1 && !(refutesAt 1 1)) = true := by decide

/-- **NC3 — the left side genuinely depends on `N`, the right side genuinely does not.**
This is the exact shape of the error: an `N`-dependent quantity equated to an `N`-free
one.  All six values of `⟨Tr H^2⟩` for `N = 1,…,6` are pairwise distinct while `C_1` is a
single number. -/
theorem nc3_lhs_N_dependent :
    (Nat.beq (wickTrace 1 1) 1 && Nat.beq (wickTrace 1 2) 4
      && Nat.beq (wickTrace 1 3) 9 && Nat.beq (wickTrace 1 4) 16
      && Nat.beq (wickTrace 1 5) 25 && Nat.beq (wickTrace 1 6) 36
      && Nat.beq (catalan 1) 1) = true := by decide

/-- **NC4 — Catalan is really Catalan, and the division in `catalan` is exact.**  Two
independent formulas agree on `k = 0,…,8`, and the values are the familiar
`1, 1, 2, 5, 14, 42, 132, 429, 1430`. -/
theorem nc4_catalan_two_ways :
    (allB (fun k => Nat.beq (catalan k) (catalanAlt k)) [0, 1, 2, 3, 4, 5, 6, 7, 8]
      && Nat.beq (catalan 0) 1 && Nat.beq (catalan 1) 1 && Nat.beq (catalan 2) 2
      && Nat.beq (catalan 3) 5 && Nat.beq (catalan 4) 14 && Nat.beq (catalan 5) 42
      && Nat.beq (catalan 6) 132 && Nat.beq (catalan 7) 429
      && Nat.beq (catalan 8) 1430) = true := by decide

/-- **NC5 — the component counter has teeth.**  At `k = 1` the single pairing leaves both
index variables free (2 classes, hence `N^2`, exactly the proof's own `Σ_{i,j} δ_ii δ_jj`).
At `k = 2` the three pairings leave `3, 1, 3` classes: the CROSSING pairing `{0,2},{1,3}`
collapses all four index variables into ONE class, which is precisely where the `+N`
of `2N^3 + N` comes from.  A counter stuck at a constant would fail here. -/
theorem nc5_components_have_teeth :
    (Nat.beq (compOfMatching 2 [(0, 1)]) 2
      && Nat.beq (compOfMatching 4 [(0, 1), (2, 3)]) 3
      && Nat.beq (compOfMatching 4 [(0, 2), (1, 3)]) 1
      && Nat.beq (compOfMatching 4 [(0, 3), (1, 2)]) 3) = true := by decide

/-- **NC6 — the pairing enumeration is complete and duplicate-free**: its cardinality is
the double factorial `(2k-1)!! = 1, 3, 15, 105, 945` for `k = 1,…,5`, it is EMPTY on an
odd number of factors, and the fuel is not exhausted (fuel `0` would give `0` matchings
and this test would fail). -/
theorem nc6_matching_counts :
    (Nat.beq (llen (matchingsOf 2)) 1
      && Nat.beq (llen (matchingsOf 4)) 3
      && Nat.beq (llen (matchingsOf 6)) 15
      && Nat.beq (llen (matchingsOf 8)) 105
      && Nat.beq (llen (matchingsOf 10)) 945
      && Nat.beq (llen (matchingsOf 3)) 0
      && Nat.beq (llen (matchingsOf 5)) 0) = true := by decide

/-- **NC7 — the two independent methods agree.**  Method A (union–find, `N^c`) and
Method B (literal enumeration of all `N^{2k}` index tuples) return the same integer on
every pair cheap enough to run both.  A bug in the union–find could not survive this. -/
theorem nc7_methods_agree :
    (bruteAgrees 1 2 && bruteAgrees 1 4 && bruteAgrees 1 6
      && bruteAgrees 2 2 && bruteAgrees 2 3 && bruteAgrees 2 4
      && bruteAgrees 3 2) = true := by decide

/-- **NC8 — the index-space enumeration used by Method B is itself complete**: `N^{2k}`
tuples, each of length `2k`. -/
theorem nc8_tuple_space_complete :
    (Nat.beq (llen (tuplesOf 2 4)) 16
      && Nat.beq (llen (tuplesOf 3 4)) 81
      && Nat.beq (llen (tuplesOf 2 6)) 64
      && allB (fun v => Nat.beq (llen v) 4) (tuplesOf 3 4)) = true := by decide

/-- **NC9 — the delta constraints are the propagator's, not something weaker.**  At
`k = 1` the pairing `{0,1}` yields the constraints `i_0 = i_0`, `i_1 = i_1` (the two
deltas `δ_ii δ_jj` the proof writes down), while at `k = 2` the crossing pairing yields
four constraints that chain all four variables together.  A `consOf` that emitted no
constraints would make every count `N^{2k}` and NC1 would fail; one that emitted too many
would collapse everything and NC1 would fail too. -/
theorem nc9_constraints_are_the_propagators :
    (Nat.beq (llen (consOf 2 [(0, 1)])) 2
      && Nat.beq (llen (consOf 4 [(0, 2), (1, 3)])) 4
      && Nat.beq (nextIdx 2 1) 0 && Nat.beq (nextIdx 4 3) 0
      && Nat.beq (nextIdx 4 1) 2) = true := by decide

/-- **NC10 — the gap is not a rounding artefact of size 1**: at `(k,N) = (3,2)` the exact
left side is 24 TIMES the asserted right side, and at `(k,N) = (5,6)` it is 60030 times
it.  Meanwhile the CORRECT asymptotic leading term `C_k · N^{k+1}` is within a factor 2 of
the exact value at `(5,6)` — the claim's error is the missing `N^{k+1}`, exactly as the
record states. -/
theorem nc10_gap_is_structural :
    (Nat.beq (wickTrace 3 2) (24 * catalan 3)
      && Nat.beq (wickTrace 5 6) 2521260
      && Nat.beq 2521260 (60030 * catalan 5)
      && Nat.ble (catalan 5 * npow 6 6) 2521260
      && Nat.ble 2521260 (2 * (catalan 5 * npow 6 6))) = true := by decide

end MSE_811862

\end{Vlean}
\subsection*{Proof-safety information}
The automated scan found no unfinished proofs, new axiom declarations, unsafe declarations, external code hooks, or uses of Lean's compiled decision procedure. This scan concerns proof construction only; it does not establish that the theorem matches the source claim.
\textit{Technical scan counts:} native\_decide=0, sorry=0, axiom\_decl=0, unsafe=0, extern=0. \textit{Technical tier:} KERNEL-ONLY.
\textbf{A YES means:} Lean represents the same objects and assumptions, uses the same proposed counterexample, and the final theorem statement asserts the same contradiction. Successful compilation by itself is not enough.
\clearpage
\section*{Human sign-off}
Complete this block from your own review of Questions 1--3. Do not use the automated appendix below as an answer key.\par
\textbf{Reviewer name:} \rule{9cm}{0.4pt}\par
\textbf{Date:} \rule{4cm}{0.4pt}\par
\medskip\textbf{1. Does the proof verification document accurately represent the source and challenged claim?}\par
\answerchoices\par
\decisionexplanation
\medskip\textbf{2. Does the proposed counterexample satisfy the assumptions and make that claim false?}\par
\answerchoices\par
\decisionexplanation
\medskip\textbf{3. Does the final Lean theorem statement represent and verify that same counterexample?}\par
\answerchoices\par
\decisionexplanation
\medskip\textbf{Overall decision:}\par
\checkbox\ VALIDATED (all three answers are YES)\par
\checkbox\ REJECTED (at least one answer is NO)\par
\checkbox\ NEEDS ANOTHER REVIEWER (at least one CANNOT DETERMINE and no NO)\par
\medskip\textbf{Additional comments.} Use this space for any extra source, counterexample, calculation, or Lean details that support your decisions.\par\rule{\linewidth}{0.4pt}\par\vspace{7pt}\rule{\linewidth}{0.4pt}
\clearpage
\section*{Appendix --- Prior automated review (background only)}
\textbf{These results are not an answer key.} They come from the earlier Opus/Fable verification campaign and are printed only so an auditor can see what was previously checked. The human decisions above must be made independently. Raw status codes are retained in small type for traceability.
\textbf{Overall automated disposition:} Cleared by automated screening; an independent search also found a counterexample.
\par\textit{Audit code:} \nolinkurl{PASS_CORROBORATED}
\subsection*{Independent search with no access to the stored example}
\textbf{Result:} An independent search found a counterexample without seeing the stored example.
\par\textit{Audit code:} \nolinkurl{WITNESS_FOUND}
\par\textbf{Checker's explanation:}
\begin{Vsrc}
The quoted line asserts a chain of two equalities and both are false as written. (1) The middle term sum_{i_1..i_{2k}} H_{i_1 i_2} H_{i_2 i_3}... H_{i_{2k} i_1} is, by the definition of matrix multiplication and trace, identically Tr(H^{2k}) for the individual matrix H -- the ensemble average has been silently dropped, so equating it to <Tr(H^{2k})> is a type error, and equating it to the constant C_k is refuted by H=[[2]] (N=1), where it equals 4 and C_1=1. (Shortened here; the complete note remains in the machine-check JSON.)
\end{Vsrc}
\par\textbf{Cases checked:}
\begin{Vsrc}
PART A (second equality as a deterministic identity): all symmetric N x N integer matrices with N in {1,2} and entries in {-2,-1,0,1,2} (5^1 + 5^3 = 130 matrices), each tested for k in {1,2,3} => 390 (matrix,k) cells. (Shortened here; the complete note remains in the machine-check JSON.)
\end{Vsrc}
\subsection*{Independent re-check of the proposed counterexample}
\textbf{Result:} A separate checker recomputed the proposed counterexample and confirmed it.
\par\textit{Audit code:} \nolinkurl{CONFIRMED}
\par\textbf{Checker's note:}
\begin{Vsrc}
Checked from scratch without reading the recorded search output or the Lean certificate. Runtime ~2 s.
\end{Vsrc}
\subsection*{Independent review of the source-to-Lean correspondence}
\textbf{Result:} A separate reviewer judged that Lean proves the same failure using the same example.
\par\textit{Audit code:} \nolinkurl{MATCHES}
\par\textbf{Theorem summary from that reviewer:}
\begin{Vsrc}
The Lean file defines wickTrace k N = the exact Wick/Isserlis evaluation of <Tr(H^{2k})> for an N x N matrix under the answer's own propagator <H_ij H_kl> = delta_il delta_jk (sum over all (2k-1)!! perfect matchings of the 2k cyclic factors of N^{#index-classes}, computed in Nat by union-find), and catalan k = binom(2k,k)/(k+1) = the claim's own right-hand side. claimHolds k N := (wickTrace k N == catalan k) is the refuted sentence instantiated at (k,N); (Shortened here; the complete note remains in the machine-check JSON.)
\end{Vsrc}
\par\textbf{Difference noted by that reviewer:}
\begin{Vsrc}
None material; all three links hold. (1) Definitions. Tr(H^{2k}) is expanded exactly as in the quoted sentence, sum over i_1..i_{2k} of H_{i_1 i_2}...H_{i_{2k} i_1}: factor t is H_{a_t b_t} with a_t = t, b_t = t+1 mod 2k (nextIdx). consOf turns the pairing {s,t} into precisely the two deltas the answer's propagator gives, delta_{a_s b_t} and delta_{b_s a_t}, i.e. the constraints (s, t+1) and (s+1, t); this is symmetric under s<->t, as it must be. (Shortened here; the complete note remains in the machine-check JSON.)
\end{Vsrc}
\subsection*{Follow-up review}
No follow-up review was needed or recorded.
\textit{This appendix is only background. You do not need any separate machine-output files to answer the three questions.}
\end{document}